\documentclass[a4paper,11pt]{article}
\usepackage[utf8]{inputenc}
\usepackage[T1]{fontenc}
\usepackage[turkish]{babel}
\usepackage[margin=1in]{geometry}
\usepackage{multicol}
\usepackage{enumitem}
\usepackage{parskip}
\usepackage{hyperref}
\usepackage{array}

\setlist[itemize]{noitemsep, topsep=0pt, leftmargin=2em}
\setlength{\columnsep}{0.5cm}

\hypersetup{colorlinks=true, urlcolor=black}

\begin{document}

\vspace*{-6.5em}
\begin{center}
    \href{https://menasy.me/}{\LARGE \textbf{Mehmet Nasim Yılmaz}}
\end{center}
\vspace*{0.5em}
\noindent
\begin{minipage}[t]{0.48\textwidth}
    \begin{tabular}{@{}l}
        \textbf{LinkedIn:} \href{https://linkedin.com/in/menasy}{linkedin.com/in/menasy} \\
        \textbf{GitHub:} \href{https://github.com/menasy}{github.com/menasy} \\
        \textbf{Portfolio:} \href{https://menasy.me/}{menasy.me} \\
    \end{tabular}
\end{minipage}%
\hfill
\begin{minipage}[t]{0.48\textwidth}
    \raggedleft
    \begin{tabular}{@{}l}
        \textbf{Lokasyon:} İstanbul \\
        \textbf{Doğum Tarihi:} 02.02.2000 \\
        \textbf{E-posta:} mehmetnasim42@gmail.com \\
    \end{tabular}
\end{minipage}

\vspace{0.5em}

% Özet
Yazılım geliştirme; web teknolojileri ve mobil programlama alanlarına odaklanıyorum. \textit{"Bilmediğim hiçbir şey yok, sadece öğrenmediğim şeyler var."} anlayışıyla kendimi sürekli geliştirmeye devam ediyorum. Bugüne kadar edindiğim teknik bilgi ve deneyimle sağlam bir altyapı oluşturduğuma inanıyorum. Şimdi ise hedefim; sahip olduğum yetkinliklerle yenilikçi projelere imza atmak, dahil olduğum ekiplere değer katmak ve sektörde iz bırakan işler başarmaktır.
\vspace{-0.2em}

Blog yazılarım için \textbf{\href{https://medium.com/@menasy}{Medium}}, tüm projelerim ve detaylı bilgiler için \textbf{\href{https://menasy.me/}{portföyüme}} göz atabilirsiniz. Şirketinizde çalışmak benim için büyük bir fırsat olacaktır. İlginiz için teşekkür eder, saygılarımı sunarım.

\vspace{-0.1em}

\textbf{Eğitim Bilgisi}
\vspace{-0.5em}
\begin{itemize}
  \item İstanbul Gelişim Üniversitesi – Yönetim Bilişim Sistemleri / 2021 - 2025
\end{itemize}

\vspace{-0.4em}

\textbf{Deneyimler}
\vspace{-0.5em}
\begin{itemize}

  \item \textbf{42 İstanbul – Yazılım Mühendisi / 09.10.2023 - 15.01.2026}
    
    C/C++, multi-threading, socket ve mikroservis mimarisinden TypeScript ile full-stack web geliştirmeye kadar çeşitli projelerde yazılım geliştirici olarak yer aldım.
  \item \textbf{Baykar Teknoloji – Merkezi Kontrol Yazılımları Stajyeri / 06.10.2025 - 27.12.2025}
  
    Wowza, .NET ve React kullanarak full-stack bir canlı yayın platformu geliştirdim.
  
  \item \textbf{Paka Teknoloji – Sistem Destek Stajyeri / 12.08.2024 – 23.09.2024}
  
  IT altyapı yönetimi ve donanım/yazılım destek süreçleri operasyonlarını gerçekleştirdim.
  
  \item \textbf{Tencent Holdings Limited – Sunucu Teknisyeni / 01.04.2023 – 01.05.2023}
  
Ekip halinde sunucuların kurulumu, rack yerleşimi ve kablolama işlemlerini tamamladık.
\end{itemize}

\vspace{-0.4em}

\textbf{Öne Çıkan Projeler}
\vspace{-0.5em}
\begin{itemize}
    \item \textbf{eSIMeco - Web:} Next.js ile geliştirdiğim, production-ready eSIM satın alma platformu.
    \item \textbf{Ft-transcendence - Web Oyunu:} Klasik Pong oyununu dört kişilik bir ekip olarak web tabanlı modern teknolojilerle baştan geliştirdik.
    \item \textbf{Webserv - Web Sunucusu:} Ekip arkadaşımla birlikte C++ ve soket programlama ile HTTP/1.1 destekleyen, NGINX işleyişine sahip bir web sunucusunu sıfırdan geliştirdik.
  \item \textbf{Inception - Mikroservis:} Docker ile web tabanlı mikroservis mimarisi oluşturdum.
    \item \textbf{HesKit - Mobil:} Java/SQL kullanarak  çalışan yönetim sistemi uygulaması yaptım.
    \item \textbf{Merkezi Sağlık Sistemi - Mobil:} Kotlin ve Firebase tabanlı randevu/reçete sistemi.
  \item \textbf{Akıllı Rehber - Mobil:} React Native ile erişilebilirlik odaklı sade ve kullanıcı odaklı özelliklere sahip rehber uygulamasını geliştirip Play Store'da yayınladım.
\end{itemize}

\vspace{-0.5em}

% İlk satır: Programlama Dilleri & Programlama Becerileri
\noindent
\begin{minipage}[t]{0.48\textwidth}
    \textbf{Programlama Dilleri}
    \vspace{0em}
    \begin{itemize}
        \item C (uzman)
        \item Git/GitHub(uzman)
        \item C\#/ASP.NET (ileri)
        \item JavaScript/TypeScript (ileri)
        \item React/Next.js/React Native (ileri)
        \item C++ (orta)
        \item Java (orta)
        \item Kotlin (orta)
        \item SQL/PostgreSQL (orta)
        \item HTML/Tailwind/CSS (orta)
        \item XML (orta)
        \item Python (temel)
    \end{itemize}
\end{minipage}%
\hfill
\begin{minipage}[t]{0.48\textwidth}
    \textbf{Programlama Becerileri}
    \vspace{0em}
\begin{itemize}
    \item Çoklu İşlem/Çoklu İş Parçacığı Yönetimi
    \item REST API Tasarımı ve Geliştirme
    \item Sunucu Taraflı Oluşturma (SSR)
    \item Docker ve Mikroservis Mimarisi
    \item Nesne Yönelimli Programlama
    \item Monitoring ve Log Yönetimi
    \item Mobil Uygulama Geliştirme
    \item Component-Based Mimari
    \item Onion ve MVC Mimarileri
    \item Asenkron Programlama
    \item Linux Sistem Yönetimi
    \item CI/CD
\end{itemize}
\end{minipage}

\vspace{0.1em}
\noindent
\begin{minipage}[t]{0.48\textwidth}
    \textbf{İlgi Alanları}
    \vspace{0em}
    \begin{itemize}
        \item Mobil Uygulama Geliştirme
        \item Full-Stack Geliştirme
        \item Web Teknolojileri
    \end{itemize}
\end{minipage}%
\hfill
\begin{minipage}[t]{0.48\textwidth}
    \textbf{Diller}
    \vspace{0em}
    \begin{itemize}
        \item İngilizce (B1)
    \end{itemize}
\end{minipage}

\end{document}